\section{The question and the scope of the results}\label{sec:scope}

The investigation asks whether the arithmetic of the complete
Riemann--Siegel representation can rule out the multiple real zero that
would occur at a positive de Bruijn--Newman threshold. There are two
distinct mathematical obligations: obtain a complete, accurately
normalized arithmetic expression, and prove a sign for that expression
that contradicts the necessary geometry of the proposed event.
The first obligation has been met in the local regime developed below.
The second remains open.

The organization follows this distinction. Sections~\ref{sec:geometry}
and~\ref{sec:representation} identify the event and a denominator-free
signed measurement. Section~\ref{sec:fixed} explains why the tested
fixed-height measurements do not resolve a fixed multiplicity.
Section~\ref{sec:micro} derives the microscopic approximation, its
coefficient moments, and the error after imposing the two event equations.
Sections~\ref{sec:phases}--\ref{sec:mirror} test progressively more
restrictive phase information. The final construction keeps all
coefficient magnitudes, complete unit-phase multiplicativity, actual
phases at a growing set of small primes, the actual carrier and correction,
both local event equations, and one necessary fourth-derivative condition.
It still has the sign permitted by a collision.

Throughout the asymptotic sections,
\begin{equation}\label{eq:band}
 L=\log\frac{a}{4\pi},\qquad T=\frac cL,\qquad
 \frac12\le c\le1,\qquad a\longrightarrow\infty.
\end{equation}
Constants are uniform in this interval of \(c\), unless a different
dependence is stated. A fixed disk radius, a fixed \(\beta>0\), or a fixed
derivative order may enter the constant. No bound stated for a fixed
parameter is used with that parameter growing without a separate proof.
The eventual objective concerns every \(0<T\le1/2\), every \(a>0\), and
every multiplicity \(m\ge2\). There is no premise locating a hypothetical
threshold event in \eqref{eq:band}.

The research classification follows the
\repo{result-summaries/README.md}{repository result index}, especially
\repo{result-summaries/11_newman_heat_flow_reductions_and_certificates.md}{cluster 11},
and the \repo{avoid/README.md}{avoidance index}, particularly
\repo{avoid/01_restating_the_target.md}{target restatement},
\repo{avoid/02_shape_only_arguments_and_the_controls.md}{shared control premises},
\repo{avoid/07_scoped_exclusions_read_as_universal.md}{scoped exclusions},
and \repo{avoid/12_term_by_term_matching_and_separate_bounds.md}{separate bounds}.
The endpoint reduction is stated once below. The measurable contributions
are the explicit acquisition errors, moment asymptotics, matrix scale,
and signed counterexamples on precisely defined enlarged classes.

\section{The canonical flow and the finite event}\label{sec:geometry}

\subsection{Normalization and the positive-threshold reduction}

We use the normalization
\begin{align}
 \Phi(u)&=\sum_{n\ge1}
 (2\pi^2n^4\ee^{9u}-3\pi n^2\ee^{5u})
 \ee^{-\pi n^2\ee^{4u}},\label{eq:source}\\
 H_t(z)&=\int_0^\infty \ee^{tu^2}\Phi(u)\cos(zu)\dd u,\qquad
 H_0(z)=\frac18\xi\!\left(\frac{1+\ii z}{2}\right).
 \label{eq:flow}
\end{align}
The rapid source decay permits differentiation on compact sets:
\begin{equation}\label{eq:heat}
 \partial_t H_t=-\partial_{zz}H_t.
\end{equation}
For real \(t\), \(H_t\) is even, real on the real axis, and entire of
order one. Positivity of \(\Phi\) on \(u\ge0\) gives
\(H_t(\ii y)>0\) for real \(y\). The finite de Bruijn--Newman threshold
\(-\infty<\Lambda\le1/2\) satisfies
\begin{equation}\label{eq:threshold}
 \{\text{all zeros of }H_t\text{ are real}\} \quad\Longleftrightarrow\quad
 t\ge\Lambda.
\end{equation}
These classical inputs use the normalization of \cite{polymath}.

\begin{proposition}[Inherited finite-attainment reduction]\label{prop:finite}
If \(\Lambda>0\), then at \(T=\Lambda\) there are a finite \(a>0\)
and an integer \(m\ge2\) such that \(a\) and \(-a\) are zeros of
\(H_T\), each of multiplicity \(m\).
\end{proposition}
\begin{proof}
Take \(t_j\uparrow\Lambda\), with \(t_j\ge\Lambda/2>0\), and choose a
nonreal zero \(z_j\) of \(H_{t_j}\). The positive-time zero-strip theorem
places these zeros in \(|\Ima z|\le1\). The uniform high-zero theorem
\cite[Theorem 1.5(i)]{polymath} makes every zero with
\(|\Rea z|\ge\exp(C/t_j)\) real. Thus, after using the symmetries and
passing to a subsequence, the \(z_j\) lie in one compact set and converge
to a zero \(a\) of \(H_\Lambda\). All zeros at \(\Lambda\) are real, and
\(H_\Lambda(0)>0\), so choose the symmetry with \(a>0\).

If \(a\) were simple, the implicit-function theorem would give a unique
nearby zero for real times near \(\Lambda\). Conjugation would make this
branch real. This contradicts the approaching nonreal zeros. Evenness
supplies the mirror zero with the same multiplicity. This is the
finite-attainment argument already recorded in \cite{collision}.
\end{proof}

For clarity, the endpoint equivalence used by the program is
\begin{equation}\label{eq:endpoint}
 \mathrm{RH}\quad\Longleftrightarrow\quad
 \text{there is no real }a\text{ with }H_t(a)=H_t'(a)=0
 \text{ for any }t>0.
\end{equation}
The reverse implication follows from Proposition~\ref{prop:finite}.
For the forward implication, all times \(t\ge0\) are all-real under RH.
A multiple real zero at \(t>0\) would split nonreally immediately
backwards. Indeed, applying the heat equation to the first nonzero
Taylor term \(b(z-a)^m\) gives
\[
 \ee^{-\tau\partial_z^2}(z-a)^m
 =(2\tau)^{m/2}
 \operatorname{He}_m\!\left(\frac{z-a}{\sqrt{2\tau}}\right).
\]
The simple real zeros of the probabilists' Hermite polynomial give local
branches \(a+\sqrt{2\tau}\,\lambda_j+O(|\tau|)\); for \(\tau<0\),
some are nonreal. This argument, including arbitrary multiplicity, is
inherited from \cite{collision,polymath}. Equation~\eqref{eq:endpoint}
is a target-equivalent reformulation, not an additional estimate.

\subsection{The reflected Poisson measurement}

Fix an all-real slice \(H_T\) and write its even product as
\[
 H_T(z)=H_T(0)\prod_{r>0}\left(1-\frac{z^2}{r^2}\right),
\]
with positive zeros repeated by multiplicity. Pairing the logarithmic
derivative, and then taking its imaginary part, gives an absolutely
convergent sum
\begin{equation}\label{eq:poisson}
 Q(a,y):=-y\Ima\frac{H_T'(a+\ii y)}{H_T(a+\ii y)}
 =\sum_r\frac{y^2}{(a-r)^2+y^2},\qquad y>0.
\end{equation}
Here and below an unqualified sum over \(r\) includes both signs and
multiplicities. The order-one product and its differentiated form are
the inputs used in \cite[Section 3]{polymath}.

Suppose \(a\) has multiplicity \(m\ge2\). Define
\begin{equation}\label{eq:ks}
 k(a,y)=1+\frac{y^2}{4a^2+y^2},\qquad
 S(a,y)=\sum_{r\notin\{a,-a\}}\frac1{(a-r)^2+y^2},
\end{equation}
removing all \(m\) copies of each event root. Then
\begin{equation}\label{eq:eventmass}
 Q(a,y)=m k(a,y)+y^2S(a,y).
\end{equation}
The mirror term is essential even when it is asymptotically small.

\begin{proposition}[Denominator-free necessary sign]\label{prop:sign}
Set, at \(z=a+\ii y\),
\begin{equation}\label{eq:B}
 \cB(a,y)=2k(a,y)|H_T(z)|^2+
             y\Ima\!\left(H_T'(z)\overline{H_T(z)}\right).
\end{equation}
At an all-real canonical event of multiplicity \(m\ge2\),
\begin{equation}\label{eq:Bnegative}
 \cB(a,y)=\bigl((2-m)k(a,y)-y^2S(a,y)\bigr)|H_T(a+\ii y)|^2<0.
\end{equation}
\end{proposition}
\begin{proof}
Multiply \eqref{eq:eventmass} by \(|H_T|^2\) and subtract it from
\(2k|H_T|^2\). All-realness ensures \(H_T(a+\ii y)\ne0\).
The canonical field has further zeros, so \(S>0\).
\end{proof}

At a fixed all-real canonical slice, existence of a \(y>0\) with
\(\cB(a,y)>0\) is equivalent to \(a\) not being a multiple zero.
If \(H_T(a)\ne0\), then \(Q(a,y)\to0\); if \(a\) is simple, then
\(Q(a,y)\to1\). Since \(k\to1\), these cases give positivity for small
\(y\). Multiple zeros are covered by \eqref{eq:Bnegative}.
This observation identifies the strength of an unrestricted sign demand;
the research task is to obtain a concrete arithmetic estimate on a
declared domain.

\subsection{What an accurate arithmetic value would have to pay}

Let \(F=H_T(z)\), and suppose \(f,g\) approximate its complete value
and physical derivative:
\[
 |F-f|\le\epsilon_0,\qquad |F'-g|\le\epsilon_1.
\]
Expanding each product yields
\begin{equation}\label{eq:Berror}
\begin{split}
 |\cB-(2k|f|^2+y\Ima(g\bar f))|
 \le{}&2k(2|f|\epsilon_0+\epsilon_0^2)\\
 &+y\bigl(|g|\epsilon_0+|f|\epsilon_1+\epsilon_0\epsilon_1\bigr).
\end{split}
\end{equation}
Thus a positive computed expression must exceed this entire error.
No denominator floor is required. If a quotient is used instead,
\(|f|>\epsilon_0\) is necessary for the bound
\begin{equation}\label{eq:quotient}
 \left|Q+y\Ima\frac gf\right|
 \le y\,\frac{\epsilon_1+|g/f|\epsilon_0}{|f|-\epsilon_0}.
\end{equation}

The cancellation at a double zero explains the sensitivity of this test.
Write \(H_T(a+w)=b_mw^m+O(w^{m+1})\), \(b_m\ne0\), and put
\(S_0=\sum_{r\ne\pm a}(a-r)^{-2}>0\).
Isolation of the event and convergence of the squared reciprocal tail
give \(S(a,y)\to S_0\). Hence
\begin{equation}\label{eq:small-y}
 \cB(a,y)=
 \begin{cases}
 -|b_2|^2 S_0 y^6(1+o(1)),&m=2,\\
 (2-m)|b_m|^2 y^{2m}(1+o(1)),&m>2.
 \end{cases}
\end{equation}
For \(m=2\), the separate terms in \eqref{eq:B} have size \(y^4\).
In the separate absolute envelope \eqref{eq:Berror},
\(\epsilon_0=o(y^4)\) and \(\epsilon_1=o(y^3)\) suffice for an
\(o(y^6)\) error. These are fixed-event sufficient conditions for that
envelope; they are not lower bounds on actual approximation errors.
